\chapter{Estimating the transition kernel}
\label{ch:estimation}

Everything so far assumed the chain is known. This chapter removes that assumption: we observe
only noised sequences and must estimate $\kernel$ from them, never seeing a clean sample.

\section{The objective}

Write $\params$ for the kernel parameters. The observed-data log-likelihood is
\begin{equation}
\label{eq:marginal-lik}
\Lik(\params) \;=\; \sum_{\mu=1}^{\nseq} \log P_{t_\mu,\params}(x^\mu),
\qquad
P_{t,\params}(x) = \int\!\mathrm{d}a\; \Pz^\params(a)\, P(x\mid a,t) .
\end{equation}
The integral is the same marginalisation Chapter~\ref{ch:bp} made tractable --- it is the
evidence, and BP already returns it.

\begin{intuition}
This is a latent-variable problem in the textbook sense: the clean sequence $a$ is the latent
variable, the noised $x$ is observed, and the parameters live in the prior. What is unusual is
that the latent structure is a tree, so the quantities EM needs are available \emph{exactly}
rather than by approximation --- which is not the typical situation.
\end{intuition}

\section{Expectation maximisation}

EM alternates two steps. Given a current $\params^{(r)}$, form
\begin{equation}
\label{eq:qfun}
\Qfun(\params\mid\params^{(r)})
= \sum_\mu \E_{\params^{(r)}}\!\left[\log P_\params(a, x^\mu)\,\middle|\, x^\mu\right],
\end{equation}
then set $\params^{(r+1)} = \arg\max_\params \Qfun$.

\begin{proposition}[Monotone ascent]
$\Lik(\params^{(r+1)}) \ge \Lik(\params^{(r)})$.
\end{proposition}

\begin{proof}
For any $\params$, write
$\log P_\params(x) = \Qfun(\params\mid\params^{(r)}) - H(\params\mid\params^{(r)})$ with
$H(\params\mid\params^{(r)}) = \E_{\params^{(r)}}[\log P_\params(a\mid x)\mid x]$. By Gibbs'
inequality $H(\params\mid\params^{(r)}) \le H(\params^{(r)}\mid\params^{(r)})$ for every
$\params$, since the difference is a Kullback--Leibler divergence. Choosing $\params^{(r+1)}$
to increase $\Qfun$ therefore increases $\Lik$ by at least as much.
\end{proof}

\begin{intuition}
Monotonicity is the reason EM is trustworthy without a step size. It is also the sharpest
available test of an implementation: an error anywhere --- in the recursion, the accumulation,
or the maximisation --- generically breaks it. In every run reported in this project the
violation is exactly zero.
\end{intuition}

\section{The expectation step is exact}

In a general latent-variable model \eqref{eq:qfun} is intractable and one resorts to a
variational or sampled approximation, losing monotonicity with it. Here the posterior is a
tree, so Proposition~\ref{prop:bp-exact} gives the required expectations exactly.

The complete-data log-likelihood is a sum over edges, so the \emph{$\params$-dependent part} of
$\Qfun$ --- it also contains a site-one term and a likelihood term, both $\params$-free and
therefore additive constants --- depends on the data only through
\begin{empheq}[box=\fbox]{equation}
\label{eq:xi-chapter}
\Xis_{kl} \;=\; \sum_{\mu=1}^{\nseq}\sum_{i=2}^{\nsites}
P\!\left(a_{i-1}=u_k,\ a_i=u_l \,\middle|\, x^\mu\right),
\end{empheq}
the expected transition mass between grid points, accumulated from the pairwise marginals
\eqref{eq:pairwise}.

\begin{intuition}
$\Xis$ is the continuum analogue of the expected transition counts of Baum--Welch. In the
discrete case it literally counts how often the chain was inferred to move from one state to
another; here it is a mass on pairs of grid points.
\end{intuition}

\begin{pitfall}
Sufficiency holds only under four conditions: a \emph{homogeneous} kernel shared across sites
and sequences, a fixed grid, an initial law excluded from $\params$, and $\params$-free
likelihood factors. The last is what licenses observations at several noise levels to
\emph{add} into one $\Xis$. Relax any of them and \eqref{eq:xi-chapter} is no longer
sufficient.

Note also what sufficiency does and does not buy. The maximisation step never revisits the
data, so it costs $O(P\ngrid^2)$ independently of $\nseq$. But \emph{forming} $\Xis$ costs
$O(\nseq\nsites\ngrid^2)$ per iteration and $O(\ngrid^2)$ storage. ``The whole E-step is one
matrix'' is a statement about sufficiency, not about inference being cheap.
\end{pitfall}

\begin{codebox}
\coderef{src/em.py}{e\_step} accumulates \eqref{eq:xi-chapter};
\coderef{src/em.py}{ExpectedStatistics} holds it; \coderef{src/em.py}{fit\_em} is the driver
with the monotonicity trace. \coderef{src/em.py}{clean\_statistics} is the $t\to0$ limit, which
needs no BP at all --- the clean-data case Marc's original suggestion referred to.
\end{codebox}

\section{Fisher's identity: a gradient without differentiating the recursion}

There is a second route to the same information. Fisher's identity states
\begin{empheq}[box=\fbox]{equation}
\label{eq:fisher-chapter}
\nabla_\params \Lik(\params)
= \sum_\mu \E_\params\!\left[\nabla_\params \log P_\params(a,x^\mu)\,\middle|\,x^\mu\right]
= \big\langle \Xis(\params),\ \nabla_\params \log \kernel_\params\big\rangle .
\end{empheq}

\begin{proof}
$\nabla \log P(x) = \nabla P(x)/P(x)$ and
$\nabla P(x) = \int\!\mathrm{d}a\, \nabla P(a,x)
= \int\!\mathrm{d}a\, P(a,x)\nabla\log P(a,x)$; dividing by $P(x)$ turns $P(a,x)/P(x)$ into
the posterior. The second equality follows because $\log P_\params(a,x)$ depends on $\params$
only through the edge terms.
\end{proof}

\begin{intuition}
This says something practically important. BP is not a computation graph to be
backpropagated through --- it is the oracle that supplies the conditional expectation. The
exact gradient of the \emph{marginal} likelihood comes out of one forward--backward pass. Advice
received elsewhere on this project held that EM here would require moving the recursion into an
automatic-differentiation framework; that is not so, and the entire implementation is plain
array arithmetic.
\end{intuition}

\begin{pitfall}
Fisher's identity is \emph{not} EM. It licenses gradient ascent on $\Lik$; EM maximises
$\Qfun$ at each step. The two are different algorithms and neither dominates. Measured here:
on the smooth Gaussian kernel gradient ascent reaches EM's optimum to $2\times10^{-9}$ nats
while needing a tuned step size (at too large a step it diverges outright); on the Laplace
kernel it attains a \emph{higher} likelihood than EM, because there the exact maximiser of
$\Qfun$ is a lattice artefact of the discretisation.
\end{pitfall}

\begin{codebox}
\coderef{src/em.py}{q\_gradient} implements \eqref{eq:fisher-chapter};
\coderef{experiments/exp\_08\_gradient\_vs\_exact\_mstep.py}{\texttt{exp\_08}} is the
comparison. The identity
is checked against a central finite difference of the exact evidence in
\codefile{tests/test\_em\_bp.py}, agreeing to a relative $1.5\times10^{-9}$.
\end{codebox}

\section{The kernel families}

\begin{center}\small
\begin{tabular}{@{}llp{7.2cm}@{}}
\toprule
class & free parameters & maximisation step \\
\midrule
\texttt{GaussianAR1Kernel} & $(\corr,\ivar)$, 2 & closed form: belief-weighted Yule--Walker \\
\texttt{LaplaceAR1Kernel} & $(\corr,b)$, 2 & closed form, via a weighted median \\
\texttt{MixtureInnovationKernel} & $\corr$ and $C$ triples, $3C$ & inner ECM sweeps \\
\texttt{MDNKernel} & network weights & gradient, with backtracking \\
\bottomrule
\end{tabular}
\end{center}

The mixture kernel is the one used for the headline results,
\begin{equation}
\label{eq:mixture-kernel}
\kernel_\params(a'\mid a) = \sum_{c=1}^{C}\pi_c\,
\mathcal{N}\!\left(a'-\corr a;\ \nu_c,\ s_c^2\right),
\end{equation}
because it keeps the linear autoregression while leaving the innovation law essentially free.
It has never been shown the density it is asked to recover.

\paragraph{Counting the parameters.} At $C=4$: $\corr$, four weights, four locations, four
scales --- thirteen numbers, subject to $\sum_c\pi_c=1$, hence \textbf{twelve free}.

\begin{pitfall}
Two subtleties in that count. First, the zero-mean condition $\sum_c\pi_c\nu_c=0$ is
\emph{not} enforced: imposing it by reprojection is not a maximisation step and broke monotone
ascent in testing, so it is monitored as a diagnostic instead (it lands at the $10^{-3}$
level). A reader who assumed the model class of \S\ref{sec:matched-covariance} literally would
compute eleven. Second, and consequently, Proposition~\ref{prop:closure-lmmse} applies to the
\emph{true} kernel but not automatically to the estimated one, whose innovation is not exactly
centred.

Third: the mixture M-step is not the exact maximiser of $\Qfun$. The complete-data problem
carries a second latent variable --- the mixture label --- so the step is itself iterative, run
as inner conditional-maximisation sweeps. That makes the algorithm \emph{generalised} EM:
monotone ascent is preserved, exact maximisation is not. It was for some time labelled
``exact block M-step'' in the source, which it never was.
\end{pitfall}

\begin{pitfall}
The inner sweep count was fixed at four, and four sweeps is not enough to settle the
innovation \emph{shape} when the outer loop is short. Measured at $\ncomp=8$, $\nseq=256$,
$t=0.22$, $\ngrid=401$, holding the outer count at $40$ and varying only the inner budget
(the truth is $\kappa_4=3$):

\begin{center}\small
\begin{tabular}{@{}lccccc@{}}
\toprule
inner sweeps & $2$ & $4$ & $8$ & $16$ & $64$ \\
\midrule
fitted $\kappa_4$ & $0.64$ & $1.68$ & $2.06$ & $2.14$ & $2.16$ \\
$\widehat\corr$ & $0.81838$ & $0.81865$ & $0.81977$ & $0.82006$ & $0.82005$ \\
$\Lik/\text{edge}$ & $-1.25791$ & $-1.25717$ & $-1.25705$ & $-1.25705$ & $-1.25707$ \\
\bottomrule
\end{tabular}
\end{center}

\noindent
The shape moves by $28\%$ between four and sixteen sweeps while $\corr$ moves in the fourth
decimal --- the same rate asymmetry the outer loop shows, where $\corr$ is flat from about
iteration $25$ and the innovation shape needs of order a hundred, reappearing one level down.
The objective is also non-monotone in this budget: $64$ sweeps lands below
$16$. So this is a third instance of the pattern that runs through the whole project, after
EM's outer budget and Rung 4a's gradient ascent: \emph{the optimisation budget is a
regularisation knob, and more ascent is not more accuracy.}

The scope of the effect matters, and the first version of this note overstated it by
generalising from the $40$-iteration column. At $60$ outer iterations the inner budget stops
mattering, because the two loops partially substitute --- four sweeps get there too, just
slower. At $\nseq=128$ the fitted $\kappa_4$ is $3.22$ at four sweeps against $3.34$ at
sixteen; at $\nseq=512$ it is $2.97$ against $2.93$, i.e.\ non-monotone and within noise. In
the efficiency run, whose outer count is selected on validation from a grid reaching $400$,
the median selection is $120$ and $70\%$ of cells sit at $60$ or above --- so most of that run
is in the regime where this does not bite, and $30\%$ of it is not.

The default is now sixteen, which removes the interaction rather than asking a reader to
work out which regime a cell was in.
\end{pitfall}

\begin{pitfall}
The same measurements settle a resolution question that had been open. A mixture component
narrower than a grid cell can raise the quadrature likelihood without corresponding to
anything in the innovation law, so an under-resolved fit must never count as evidence for
higher capacity. At $\ncomp=8$, $\ngrid=401$ the narrowest fitted component is
$s_{\min}/\gridh \approx 9.4$--$10.0$: nearly ten cells wide, five times the threshold below
which a component is a lattice artefact. The production fits are resolved with room to spare.
This is now recorded per cell by the efficiency experiment rather than spot-checked, and
\codefile{tools/make\_tab\_efficiency.py} refuses to build the table from cells that fail it.
\end{pitfall}

\begin{figure}[t]
\centering
\includegraphics[width=0.98\textwidth]{fig_em_diagnostics.pdf}
\caption{Left: the marginal log-likelihood by iteration at $\ncomp=4$, $\nseq=512$, from four
initialisations of $\corr$; the largest decrease across iterations is zero in all four runs, and
in the four $\ncomp=8$ runs not shown. Centre: the estimated autoregressive coefficient over the
same runs, reaching $[0.8517,0.8520]$ against a truth of $0.85$ from initialisations spanning
$[-0.4,0.6]$ --- which is the direct evidence that $\corr$ is estimated rather than supplied.
(The eight $\ncomp\in\{4,8\}$ runs together span $[0.8517,0.8529]$.)
Right: RMSE of $\widehat\corr$ and of the innovation variance against the number of training
sequences, with an $\nseq^{-1/2}$ reference. Sources \codefile{outputs/exp\_18/em\_trace.csv},
\codefile{outputs/exp\_06\_em\_parameter\_recovery/em\_rate.csv}.}
\label{fig:em-compendium}
\end{figure}


\begin{codebox}
\coderef{src/kernels.py}{MixtureInnovationKernel} --- see its \texttt{m\_step}, whose
\texttt{n\_inner} is the ECM sweep cap (now $16$) and \texttt{inner\_tol} the $\Qfun$-based
inner criterion; the returned kernel carries \texttt{inner\_sweeps},
\texttt{inner\_converged} and \texttt{inner\_q\_gain} so a starved inner loop is visible.
\texttt{theta} exposes the flat parameter vector, which is what the parameter count is read
from: $1+3\ncomp$ numbers with $\pi$ simplex-constrained, hence $3\ncomp$ free --- twelve at
$\ncomp=4$, and \textbf{twenty-four} at the $\ncomp=8$ the frozen config actually uses. The
efficiency table said twelve for longer than it should have.
\end{codebox}

\section{Identifiability}

Two claims must be separated, and confusing them overstates what is proved.

\paragraph{Distribution level.} Gaussian convolution at finite $t$ is injective. In
characteristic-function terms
$\varphi_{t,\params}(\xi) = \varphi_\params(e^{-t}\xi)\,e^{-\Dt\|\xi\|^2/2}$; the Gaussian
factor has no zeros so it divides out, and $e^{-t}>0$ makes $\xi\mapsto e^{-t}\xi$ a bijection.
So $\Pt$ determines $\Pz$.

\paragraph{Parameter level.} That $\Pz$ is determined does \emph{not} by itself determine
$\kernel$. That additionally requires: the initial law fixed or separately identifiable; the
transition family identifiable as a parametric family; the mixture treated modulo label
permutation; the mixture parameters in the interior --- all $\pi_c>0$ with the $(\nu_c,s_c)$
distinct, which an over-specified $C$ need not respect; and $\mu$ of full support, since
$\kernel$ is pinned down only there. We assume these rather than prove them.

\begin{intuition}
Injective is not the same as stable. This is Gaussian deconvolution, whose inverse is
unbounded, so identifiability survives at any finite $t$ while the \emph{information} degrades
rapidly --- and unevenly. Measuring that unevenness needs care about what is being compared. In
the Gaussian model with parameters $(\rho, q)$ the innovation \emph{scale} $q$ falls by a factor
$142$ between $t=0.05$ and $t=1.6$ while the correlation falls by $26$, but both of those are
second-order quantities: $q$ is a scale, not a shape, so the ratio says nothing about
higher-order structure.

A genuine shape coordinate is needed. On a generalised-Gaussian innovation
$\varphi_\beta(e)\propto\exp(-(|e|/s)^\beta)$ with $s$ chosen so the variance is exactly $q$,
$\beta$ moves the fourth moment at fixed second moment. Reporting the \emph{efficient}
information for each coordinate --- projecting out the other two, since the diagonal ignores
their coupling --- the decay across the same range is $56$--$87$ for the correlation,
$235$--$522$ for the variance and $8.7\times10^3$--$6.8\times10^4$ for the shape. Higher-order
structure is therefore not five times more fragile but $112$ to $1222$ times, which is exactly
the structure this project cares about.
\end{intuition}

\begin{codebox}
\coderef{experiments/exp\_06\_em\_parameter\_recovery.py}{} measures the information decay
(\texttt{price\_of\_noising.csv}) and the recovery of an unknown innovation law from noisy data
alone.
\end{codebox}

\section*{Sources}
\addcontentsline{toc}{section}{Sources}
EM is \citet{dempster1977}; the convergence analysis, including what monotone ascent does and
does not guarantee, is \citet{wu1983}. The generalised and conditional-maximisation variants
used here --- which is what the mixture M-step actually is --- are \citet{meng1993}. Fisher's
identity for the gradient of the marginal likelihood goes back to \citet{fisher1925} and is
given in the latent-variable form used here by \citet{louis1982} and \citet{cappe2005}. The
transition-count statistic $\Xis$ is the continuum analogue of the one in \citet{baum1970}.
Identifiability of finite mixtures, which Assumption~1 of the note leans on, is
\citet{teicher1963} and \citet{yakowitz1968}.
